% eksperymentalne tłumaczenie na włoski

%

\chapter{trygonometria} % lang-pl

\chapter{trigonometria} % lang-it
\index{trygonometria}
\index{trigonometria} % lang-it

% \section{Jednokładność} Podobieństwo figur, trójkątów (cechy), stosunek pól figur podobnych.

Trygonometria (od greckiego τρίγωνον (trójkąt) oraz μέτρον (miara)) to dział geometrii, który bada relacje między długościami boków trójkąta, a miarami jego kątów. % lang-pl
La trigonometria (dal greco τρίγωνον, triangolo, e μέτρον, misura) è il ramo della geometria che studia le relazioni tra le lunghezze dei lati di un triangolo e le ampiezze dei suoi angoli. % lang-it


O funkcjach trygonometrycznych kąta ostrego pisze Guzicki \cite[s. 240-254]{guzicki_2021}, Bogdańska, Neugebauer \cite[s. 88]{neugebauer_2018}. % lang-pl
Delle funzioni trigonometriche di un angolo acuto trattano Guzicki \cite[pp.~240--254]{guzicki_2021} e Bogdańska, Neugebauer \cite[p.~88]{neugebauer_2018}. % lang-it

% https://en.wikipedia.org/wiki/List_of_trigonometric_identities#Identities_without_variables
% https://en.wikipedia.org/wiki/Kunstweg

%

\pol{\section{Historia}}
\ita{\section{Storia}}

\pol{Nie będziemy zanudzać Czytelnika sto sześćdziesięcioma trzema wzorami redukcyjnymi, na funkcje i kofunkcje, na sumy i różnice kątów, blablabla.}
\ita{Non annoieremo il lettore con centosessantatré formule di riduzione, per funzioni e cofunzioni, per somme e differenze di angoli, e così via.}
\pol{Jest pełno innych książek, które to robią.}
\ita{Esistono già molti altri libri che lo fanno.}
\pol{Opowiemy trochę o etymologii różnych słów i tym, jak trygonometria rozwijać się będzie na przestrzeni wieków, co powinno być dużo ciekawsze.}
\ita{Racconteremo invece qualcosa sull'etimologia di vari termini e su come la trigonometria si sia sviluppata nel corso dei secoli, il che dovrebbe essere molto più interessante.}

\pol{Współczesne ,,sinus'' i ,,kosinus'' wywodzą się z błędu w tłumaczeniu!}
\ita{I moderni termini ``seno'' e ``coseno'' derivano da un errore di traduzione!}
\pol{Na początku będzie sanskryt i~\emph{jyā} (cięciwa łuku, ze starogreckiego χορδή -- sznurek, struna), a potem transkrypcja do arabskiego \emph{jība}, gdzie pomija się samogłoski (!) i prowadzi do jb, słowa bez żadnego znaczenia.}
\ita{All'inizio vi è il sanscrito \emph{jyā} (corda di un arco; dal greco antico χορδή, cioè corda o stringa), poi la trascrizione araba \emph{jība}, nella quale si omettono le vocali (!) e si ottiene jb, una parola priva di significato.}
% (جب)
\pol{Jb zinterpretują jako \emph{jayb}, które już coś znaczy (po angielsku: ,,bosom, pocket, fold''); Gerard z~Cremony przetłumaczy to do łacińskiego \emph{sinus} (czyli ,,zatoki'', ale też ,,zwisającego fałdu togi na piersi'')...}
\ita{Jb verrà interpretato come \emph{jayb}, che invece ha un significato (in inglese: ``bosom, pocket, fold''); Gerardo da Cremona lo tradurrà in latino come \emph{sinus} (cioè ``baia'', ma anche ``piega della toga sul petto'')...}
\pol{Kosinus to skrót od \emph{complementi sinus} jak w \emph{Canon triangulorum} Edmunda Guntera z 1620 roku; po drodze będzie też Fibonacci oraz jego \emph{sinus rectus arcus}, co pomoże w przyjęciu się ostatecznej wersji.}
\ita{Coseno è un'abbreviazione di \emph{complementi sinus}, come nel \emph{Canon triangulorum} di Edmund Gunter del 1620; nel mezzo vi saranno anche Fibonacci e il suo \emph{sinus rectus arcus}, che contribuiranno all'affermazione della forma definitiva.}

\pol{Nazwy dwóch kolejnych funkcji trygonometrycznych także mają korzenie w łacinie, bowiem \emph{tangens} znaczy ,,dotykający'', zaś \emph{secans} to ,,przecinający''.}
\ita{Anche i nomi di altre due funzioni trigonometriche hanno radici latine: \emph{tangens} significa ``che tocca'', mentre \emph{secans} significa ``che taglia''.}
\pol{Ma to sens, jeśli popatrzymy na okrąg jednostkowy.}
\ita{Ciò ha senso se si osserva il cerchio unitario.}
\pol{Wreszcie termin ,,trygonometria'' bierze się z greckiego τρίγωνον (trigonon -- trójkąt) oraz μέτρον (metron -- miara).}
\ita{Infine, il termine ``trigonometria'' deriva dal greco τρίγωνον (\emph{trigonon}, triangolo) e μέτρον (\emph{metron}, misura).}
% TODO: https://en.wikipedia.org/wiki/History_of_trigonometry#Ancient_Near_East
% TODO: https://en.wikipedia.org/wiki/History_of_trigonometry#Classical_antiquity
% TODO: https://en.wikipedia.org/wiki/History_of_trigonometry#Indian_mathematics
% TODO: https://en.wikipedia.org/wiki/History_of_trigonometry#Chinese_mathematics
% TODO: https://en.wikipedia.org/wiki/History_of_trigonometry#Islamic_world
% TODO: https://en.wikipedia.org/wiki/History_of_trigonometry#European_renaissance
% TODO: https://en.wikipedia.org/wiki/History_of_trigonometry#Modern

\pol{Podstawy trygonometrii rozwinie Hipparchos z Nikei: około 140 r.p.n.e opracuje tablicę cięciw, czyli odpowiednik współczesnych sinusów.}
\ita{Le basi della trigonometria furono sviluppate da Ipparco di Nicea: intorno al 140 a.C. compilò una tavola delle corde, equivalente agli attuali seni.}
\pol{Pierwsza księga poświęcona temu niezależnie od astronomii, \emph{,,Kitāb al-Shakl al-qattā''} (\emph{Rozprawa o czworobokach}), wyjdzie spod pióra Nasira al Dina al-Tusiego około 1250 roku.}
\ita{Il primo libro dedicato all'argomento indipendentemente dall'astronomia, \emph{``Kitāb al-Shakl al-qattā''} (\emph{Trattato sui quadrilateri}), fu scritto da Nasir al-Din al-Tusi intorno al 1250.}
% https://en.wikipedia.org/wiki/Nasir_al-Din_al-Tusi#Mathematics
\index[persons]{al Din al-Tusi, Nasir}%

%
% eksperymentalne tłumaczenie na włoski

%

\section{Twierdzenie sinusów} % lang-pl
\index{twierdzenie!sinusów}% % lang-pl
\section{Teorema dei seni} % lang-it
\index{teorema!dei seni}% % lang-it

\begin{proposition}
[twierdzenie sinusów] % lang-pl
[teorema dei seni] % lang-it
    Niech trójkąt o bokach długości $a, b, c$ oraz kątach $\alpha$, $\beta$, $\gamma$ naprzeciwko tychże boków (w tej samej kolejności) będzie wpisany w okrąg o promieniu $R$. % lang-pl
    Wtedy % lang-pl
    Sia un triangolo con lati di lunghezze $a, b, c$ e angoli $\alpha$, $\beta$, $\gamma$ opposti a questi lati (nello stesso ordine), inscritto in un cerchio di raggio $R$. % lang-it
    Allora % lang-it
    \label{twierdzenie_sinusow}%
    \begin{equation}
        \frac{a}{\sin \alpha} = \frac{b}{\sin \beta} = \frac{c}{\sin \gamma} = 2R.
    \end{equation}
\end{proposition}

Ptolemeusz w II wieku będzie wiedzieć o równoważnej formie, że boki trójkąta są proporcjonalne do cięciw podwojonych kątów naprzeciwko. % lang-pl
Brahmagupta w VII wieku wyrazi promień okręgu opisanego jako iloczyn dwóch boków trójkąta podzielony przez podwojoną wysokość, stąd można już wyprowadzić twierdzenie. % lang-pl
Sferyczną wersję pierwszy poda któryś z~muzułmańskich matematyków X wieku: Abu-Mahmud al-Khujandi, Abu Nasr Mansur albo Abu al-Wafa. % lang-pl
Tolomeo nel II secolo conoscerà una forma equivalente, secondo cui i lati del triangolo sono proporzionali alle corde dei doppi angoli opposti. % lang-it
Brahmagupta nel VII secolo esprimerà il raggio del cerchio circoscritto come il prodotto di due lati del triangolo diviso per il doppio dell'altezza; da qui si può già dedurre il teorema. % lang-it
Una versione sferica sarà data per primo da uno dei matematici musulmani del X secolo: Abu-Mahmud al-Khujandi, Abu Nasr Mansur oppure Abu al-Wafa. % lang-it

Coxeter \cite[s. 28, 29]{coxeter_1967} wyprowadza twierdzenie sinusów z symetrii okręgu opisanego na trójkącie. % lang-pl
Hartshorne \cite[s. 185]{hartshorne_2000} podaje je jako ćwiczenie. % lang-pl
Bogdańska, Neugebauer \cite[s. 89]{neugebauer_2018} oprą się o fakt \ref{maybe_3_28}, że kąty wpisane oparte o ten sam łuk mają równe miary. % lang-pl
Coxeter \cite[pp. 28, 29]{coxeter_1967} deduce il teorema dei seni dalla simmetria del cerchio circoscritto al triangolo. % lang-it
Hartshorne \cite[p. 185]{hartshorne_2000} lo propone come esercizio. % lang-it
Bogdańska e Neugebauer \cite[p. 89]{neugebauer_2018} si baseranno sul fatto \ref{maybe_3_28}, che gli angoli inscritti insistenti sullo stesso arco hanno misure uguali. % lang-it

%
% eksperymentalne tłumaczenie na włoski

%

\section{Twierdzenie kosinusów} % lang-pl

\index{twierdzenie!kosinusów}% % lang-pl
\section{Teorema del coseno} % lang-it
\index{teorema!del coseno}% % lang-it

Twierdzenie kosinusów jest bardzo stare. % lang-pl

Euklides rozpatrzy je w Elementach osobno dla trójkątów rozwartokątnych (II.12) i ostrokątnych (II.13). % lang-pl
Il teorema del coseno è molto antico. % lang-it

Euclide lo tratta negli Elementi separatamente per i triangoli ottusangoli (II.12) e acutangoli (II.13). % lang-it
% TODO: https://en.wikipedia.org/wiki/Law_of_cosines The cases of obtuse triangles and acute triangles (corresponding to the two cases of negative or positive cosine) are treated separately, in Propositions II.12 and II.13:[1]
Perski matematyk Jamshid al-Kashi znajdzie wartość $2\pi$ z~dokładnością do szesnastu cyfr, będzie autorem najdokładniejszych tablic trygonometrycznych swoich czasów i poda równoważną postać wzoru, % lang-pl

Il matematico persiano Jamshid al-Kashi calcolerà il valore di $2\pi$ con precisione fino a sedici cifre, sarà autore delle tavole trigonometriche più precise del suo tempo e fornirà una forma equivalente della formula, % lang-it
\index[persons]{al-Kashi, Jamshid}
\begin{equation}
	c = \sqrt{(b - a \cos \gamma)^2 + (a \sin \gamma)^2}
\end{equation}
dla ostrego kąta $\gamma$. % lang-pl

Taka sama metoda rozwiązywania trójkątów pojawi się w Europie w 1464 roku, kiedy Regiomontanus opublikuje \emph{,,De triangulis omnimodis''} (czyli \emph{,,O trójkątach wszelkiego rodzaju''}). % lang-pl
per l'angolo acuto $\gamma$. % lang-it

Lo stesso metodo per risolvere i triangoli apparirà in Europa nel 1464, quando Regiomontano pubblicherà \emph{De triangulis omnimodis} (cioè \emph{Sui triangoli di ogni tipo}). % lang-it
\index[persons]{Regiomontanus}%
Współczesna forma twierdzenia to zasługa François Viète'a. % lang-pl

La forma moderna del teorema è merito di François Viète. % lang-it
\index[persons]{Viète, François}%

\begin{proposition}[twierdzenie kosinusów] % lang-pl
	W trójkącie o bokach długości $a, b, c$ z kątem $\gamma$ naprzeciwko krawędzi $c$ zachodzi % lang-pl
\begin{proposition}[teorema del coseno] % lang-it
	
	In un triangolo con lati di lunghezza $a, b, c$ e angolo $\gamma$ opposto al lato $c$ vale % lang-it
	\label{twierdzenie_kosinusow}%
	\begin{equation}
		c^2 = a^2 + b^2 - 2ab \cos \gamma.
	\end{equation}
	% https://en.wikipedia.org/wiki/Law_of_cosines
\end{proposition}

Znamy różne dowody tego twierdzenia: korzystające z twierdzenia Pitagorasa, trzech wysokości trójkąta, twierdzenia Ptolemeusza, geometrii koła albo twierdzenia sinusów. % lang-pl

We Francji twierdzenie kosinusów do późna będzie nazywane \emph{théorème d'Al-Kashi}, na cześć wspomnianego wyżej Persa. % lang-pl
Conosciamo varie dimostrazioni di questo teorema: basate sul teorema di Pitagora, sulle tre altezze del triangolo, sul teorema di Tolomeo, sulla geometria del cerchio o sul teorema dei seni. % lang-it
In Francia il teorema del coseno sarà a lungo chiamato \emph{théorème d'Al-Kashi}, in onore del suddetto persiano. % lang-it


Piszą o nim Guzicki \cite[s. 258]{guzicki_2021}, Hartshorne \cite[s. 185]{hartshorne2000} jako ćwiczenie. % lang-pl
Ne scrivono Guzicki \cite[s. 258]{guzicki_2021}, Hartshorne \cite[s. 185]{hartshorne2000} come esercizio. % lang-it

\begin{theorem}[Stewarta, 1746] % lang-pl
\begin{theorem}[di Stewart, 1746] % lang-it
\label{twierdzenie_stewarta}%
\index{twierdzenie!Stewarta}%
	W trójkącie $\triangle ABC$ o bokach długości $a, b, c$ poprowadzono czewianę z~wierzchołka $C$ do boku $AB$ o długości $d$, dzieląc ten bok na odcinki długości $n$ oraz $m$, jak na rysunku: % lang-pl
	
	In un triangolo $\triangle ABC$ con lati di lunghezza $a, b, c$ è stata tracciata la ceviana dal vertice $C$ al lato $AB$ di lunghezza $d$, dividendo tale lato in segmenti di lunghezza $n$ e $m$, come mostrato nella figura: % lang-it
	\begin{center}
\begin{comment}
    \begin{tikzpicture}[scale=.4]
        %\tkzInit[xmin=-0.5,xmax=6.5, ymin=-0.5,ymax=4.5]
        % \tkzClip
        \tkzDefPoint(0, 0){A}
		\tkzDefPoint(3.25, 3.25){d}

		\tkzDefPoint(6, 0){AB}
        \tkzDefPoint(10, 0){B}
        \tkzDefPoint(1, 7){C}
        \tkzDefPoint(35:4.75){CC}
		\tkzDrawSegments(C,AB)
        \tkzDrawPolygon[line width=0.3mm](A,B,C)

        \tkzLabelPoint[below left](A){$A$}
        \tkzLabelPoint[below right](B){$B$}
        \tkzLabelPoint[above](C){$C$}

        \tkzLabelPoint(d){$d$}

		\tkzDrawPoints[size=3,color=black,fill=black!80](A,B,C,AB)
		\tkzDrawSegment[dim={$\,\,c\,\,$,-16pt,transform shape}](A,B)
		\tkzDrawSegment[dim={$\,\,n\,\,$,-8pt,transform shape}](A,AB)
		\tkzDrawSegment[dim={$\,\,m\,\,$,-8pt,transform shape}](AB,B)
		\tkzDrawSegment[dim={$\,\,b\,\,$,8pt,transform shape,sloped}](A,C)
		\tkzDrawSegment[dim={$\,\,a\,\,$,-8pt,transform shape,sloped}](B,C)
    \end{tikzpicture}
\end{comment}
    \end{center}
	Wtedy % lang-pl
	
	Allora % lang-it
	\begin{equation}
		b^2 m + c^2 n = a (d^2 + mn).
	\end{equation}
\end{theorem}

Twierdzenie bez dowodu opublikuje Matthew Stewart w 1746 roku, chociaż Coxeter przypuści, że mogło być znane nawet Archimedesowi. % lang-pl

Il teorema, senza dimostrazione, sarà pubblicato da Matthew Stewart nel 1746, anche se Coxeter sospetta che fosse noto persino ad Archimede. % lang-it
\index[persons]{Stewart, Matthew}%
\index[persons]{Archimedes}% % lang-pl
\index[persons]{Archimede}% % lang-it
% Coxeter, H.S.M.; Greitzer, S.L. (1967), Geometry Revisited, New Mathematical Library #19, The Mathematical Association of America, ISBN 0-88385-619-0 strona 6
Dowody podadzą też Simpson, Euler, Carnot i pewnie jeszcze ktoś. % lang-pl

Dimostrazioni saranno fornite anche da Simpson, Euler, Carnot e probabilmente altri. % lang-it
\index[persons]{Simpson, Thomas}%
\index[persons]{Euler, Leonhard}%
\index[persons]{Carnot, L.M.N.}%
W przyszłości często pokazywać się będzie je jako zastosowanie twierdzenia kosinusów, tak jak Guzicki \cite[s. 265]{guzicki_2021}, Bogdańska, Neugebauer \cite[s. 90-91]{neugebauer_2018}; inne rozwiązanie skorzysta z~odcinków skierowanych. % lang-pl

In futuro saranno spesso presentati come applicazione del teorema del coseno, come Guzicki \cite[s. 265]{guzicki_2021}, Bogdańska, Neugebauer \cite[s. 90-91]{neugebauer_2018}; un'altra soluzione utilizzerà segmenti orientati. % lang-it
% TODO: Zetel, s. 31 podaje bez kosinusów, proste rachunki
Znalezienie go to ćwiczenie podane przez Evesa \cite[s. 58]{eves1_1972}. % lang-pl

Trovarlo è un esercizio proposto da Eves \cite[s. 58]{eves1_1972}. % lang-it

\begin{corollary}[twierdzenie Apoloniusza] % lang-pl
	W trójkącie $ABC$ poprowadzono środkową $AD$. % lang-pl
\begin{corollary}[teorema di Apollonio] % lang-it
	
	Wtedy suma pól kwadratów zbudowanych na bokach $AB$ i $AC$ jest dwa razy mniejsza od sumy pól kwadratów zbudowanych na środkowej $AD$ i połowie podstawy ($BD$): % lang-pl
	In un triangolo $ABC$ è stata tracciata la mediana $AD$. % lang-it
	
	Allora la somma delle aree dei quadrati costruiti sui lati $AB$ e $AC$ è la metà della somma delle aree dei quadrati costruiti sulla mediana $AD$ e sulla metà della base ($BD$): % lang-it
	\begin{equation}
		|AB|^2 + |AC|^2 = 2|BD|^2 + 2 |AD|^2.
	\end{equation}
	% https://en.wikipedia.org/wiki/Apollonius%27s_theorem
\end{corollary}

Twierdzenie przetrwa jako fakt VII.122 w zbiorach Pappusa z Aleksandrii, choć nie brak będzie przypuszczeń, że występowało wcześniej w zaginionym tekście \emph{,,Plane Loci''} Apoloniusza z Pergi. % lang-pl

(Tekst będzie wspomniany na stronie \pageref{plane_loci}). % lang-pl
Il teorema sopravvive come fatto VII.122 nelle collezioni di Pappo di Alessandria, anche se non mancheranno supposizioni che fosse presente in un testo perduto \emph{Plane Loci} di Apollonio di Perge. % lang-it

(Il testo sarà menzionato a pagina \pageref{plane_loci}). % lang-it

\begin{corollary}
	W trójkącie $\triangle ABC$ o bokach długości $a, b, c$ poprowadzono środkową oraz dwusieczną z~wierzchołka $C$. % lang-pl
	
	Długość środkowej wynosi % lang-pl
	In un triangolo $\triangle ABC$ con lati di lunghezza $a, b, c$ è stata tracciata la mediana e la bisettrice dall'angolo $C$. % lang-it
	
	La lunghezza della mediana è % lang-it
	\begin{equation}
		\sqrt{\frac{a^2 + b^2}{2} - \frac{c^2}{4}},
	\end{equation}
	zaś dwusieczna ma długość % lang-pl
	
	mentre la bisettrice ha lunghezza % lang-it
	\begin{equation}
		\frac{\sqrt{ab (a+b+c)(a+b-c)}}{a+b}.
	\end{equation}
\end{corollary}

%
% eksperymentalne tłumaczenie na włoski

%

\section{Twierdzenie tangensów}
\index{twierdzenie!tangensów}%

Twierdzenie odkryje Abu al-Wafa Buzjani w X wieku oraz niezależnie Ibn Mu'adh al-Jayyani (arabski matematyk, autor pierwszego traktatu o trygonometrii sferycznej) w XI wieku. % lang-pl

Il teorema fu scoperto da Abu al-Wafa Buzjani nel X secolo e indipendentemente da Ibn Mu'adh al-Jayyani (matematico arabo, autore del primo trattato di trigonometria sferica) nell'XI secolo. % lang-it
\index[persons]{Buzjani, Abu al-Wafa}%
\index[persons]{Ibn Mu'adh al-Jayyani}%
% The law of tangents for spherical triangles was described in the 13th century by Persian mathematician Nasir al-Din al-Tusi (1201–1274), who also presented the law of sines for plane triangles in his five-volume work Treatise on the Quadrilateral.
Po polsku często padać będzie nazwa wzór Regiomontana. % lang-pl

In polacco viene spesso usata la denominazione formula di Regiomontano. % lang-it
\index{wzór Regiomontana}%

\begin{proposition}[twierdzenie tangensów] % lang-pl
	W trójkącie o bokach długości $a, b, c$ z kątami $\alpha$, $\beta$ naprzeciwko krawędzi $a$, $b$ zachodzi % lang-pl
\begin{proposition}[teorema delle tangenti] % lang-it
	
	In un triangolo con lati di lunghezza $a, b, c$ e angoli $\alpha$, $\beta$ opposti ai lati $a$, $b$ vale % lang-it
	\label{twierdzenie_tangensow}%
	\begin{equation}
		\frac{a-b}{a+b} = \frac{\tan \frac 1 2 (\alpha - \beta)}{\tan \frac 1 2 (\alpha + \beta)}.
	\end{equation}
\end{proposition}

Prosty dowód wykorzystuje twierdzenie sinusów. % lang-pl

Una dimostrazione semplice utilizza il teorema dei seni. % lang-it
\index{twierdzenie!sinusów}%
Inny korzysta ze wzorów Mollweidego (to jest faktu \ref{wzor_mollweidego}), które wyprowadza się i tak z twierdzenia sinusów. % lang-pl

Un altro usa le formule di Mollweide (cioè il fatto \ref{wzor_mollweidego}), che comunque si ricavano dal teorema dei seni. % lang-it
\index{wzór!Mollweidego}%

Dopóki nie zostanie wynaleziony komputer albo chociaż kalkulator elektroniczny, twierdzenie tangensów będzie mieć przewagę nad twierdzeniem kosinusów, ponieważ nie wymaga sięgania po tablice logarytmiczne do znalezienia pierwiastka % lang-pl

Finché non verranno inventati il computer o almeno una calcolatrice elettronica, il teorema delle tangenti avrà un vantaggio rispetto al teorema del coseno, poiché non richiede l’uso di tavole logaritmiche per estrarre la radice % lang-it
\index{twierdzenie!kosinusów}%
\begin{equation}
c = \sqrt{a^2 + b^2 - 2 a b \cos \gamma}.
\end{equation}

Po wynalezieniu rzeczonej maszyny, przewaga nie zniknie. % lang-pl

Twierdzenie tangensów pozwala uniknąć błędów numerycznych, kiedy $a \approx b$, zaś kąt $\gamma$ jest bardzo mały. % lang-pl
Dopo l’invenzione di tale macchina, il vantaggio non scomparirà. % lang-it
Il teorema delle tangenti permette di evitare errori numerici quando $a \approx b$ e l’angolo $\gamma$ è molto piccolo. % lang-it


Istnieje rzadko stosowane uogólnienie do czworokątów cyklicznych (wpisanych w okrąg): % lang-pl
Esiste una generalizzazione poco usata ai quadrilateri ciclici (inscritti in una circonferenza): % lang-it

% źródło: https://de.wikipedia.org/wiki/Tangenssatz#Verallgemeinerung_für_Sehnenvierecke
\begin{proposition}
	W czworokącie $ABCD$ wpisanym w okrąg, o bokach $a = AB$, $b = BC$, $c = CD$, $d = DA$ oraz kątach $\alpha = \angle DAB$, $\beta = \angle ABC$ zachodzi % lang-pl
	
	In un quadrilatero $ABCD$ inscritto in una circonferenza, con lati $a = AB$, $b = BC$, $c = CD$, $d = DA$ e angoli $\alpha = \angle DAB$, $\beta = \angle ABC$ vale % lang-it
	\begin{equation}
		\frac{(a+c)(b+d)}{(a-c)(b-d)} = \frac{\tan \frac 1 2 (\alpha + \beta)}{\tan \frac 1 2 (\alpha - \beta)}.
	\end{equation}
\end{proposition}

%
% eksperymentalne tłumaczenie na włoski

%

\section{Rozwiązywanie trójkątów} % lang-pl
\section{Risoluzione dei triangoli} % lang-it
% https://en.wikipedia.org/wiki/Solution_of_triangles

% TODO: https://en.wikipedia.org/wiki/Skinny_triangle

Rozwiązywanie trójkątów (łacińskie \emph{solutio triangulorum}) to najważniejsze rachunkowe zadanie w~trygonometrii. % lang-pl
Dane są pewne informacje o trójkącie na płaszczyźnie lub sferze, takie jak: długości boków albo kąty wewnętrzne; należy ustalić miary pozostałych. % lang-pl
Zastosowania obejmują geodezję, astronomię, nawigację i nie tylko. % lang-pl
La risoluzione dei triangoli (dal latino \emph{solutio triangulorum}) è il compito computazionale più importante in trigonometria. % lang-it
Sono dati alcuni dati su un triangolo nel piano o sulla sfera, come: lunghezze dei lati o angoli interni; occorre determinare le misure degli altri. % lang-it
Le applicazioni includono geodesia, astronomia, navigazione e non solo. % lang-it

Ograniczymy się do płaszczyzny i podamy kilka nowych zależności trygonometrycznych, które pomagają w rozwiązywaniu trójkątów. % lang-pl
Znamy już twierdzenia sinusów i kosinusów; w ogólności to drugie jest bezpieczniejsze od pierwszego (ponieważ z faktu, że $\sin \alpha = \frac 1 2$ nie wynika, czy kąt $\alpha$ jest ostry, czy rozwarty, kosinus zaś jest jednoznaczny). % lang-pl
Ci limiteremo al piano e presenteremo alcune nuove relazioni trigonometriche che aiutano nella risoluzione dei triangoli. % lang-it
Conosciamo già il teorema dei seni e dei coseni; in generale il secondo è più sicuro del primo (poiché dal fatto che $\sin \alpha = \frac 1 2$ non si deduce se l'angolo $\alpha$ è acuto o ottuso, mentre il coseno è univoco). % lang-it

W każdym z poniższych stwierdzeń mamy trójkąt o kątach $\alpha, \beta, \gamma$, bokach $a, b, c$, połowie obwodu $p = \frac 1 2 (a + b + c)$, promieniu okręgu opisanego $R$ i wpisanego $r$. % lang-pl
In ciascuna delle seguenti affermazioni consideriamo un triangolo con angoli $\alpha, \beta, \gamma$, lati $a, b, c$, semiperimetro $p = \frac 1 2 (a + b + c)$, raggio del circolo circoscritto $R$ e raggio del circolo inscritto $r$. % lang-it

\begin{proposition}
    Zachodzi % lang-pl
    Vale % lang-it
    \begin{equation}
    S = 2 R^2 \sin \alpha \sin \beta \sin \gamma.
    \end{equation}
\end{proposition}

\begin{proposition}
    Zachodzi % lang-pl
    Vale % lang-it
    \begin{equation}
        p = R (\sin \alpha + \sin \beta + \sin \theta).
    \end{equation}
\end{proposition}

\begin{proposition}
    Zachodzi % lang-pl
    Vale % lang-it
    \begin{equation}
        r = 4R \sin \frac \alpha 2 \sin \frac \beta 2 \sin \frac \gamma 2.
    \end{equation}
\end{proposition}

\begin{proposition}
[wzór Newtona] % lang-pl
[formula di Newton] % lang-it
\index{wzór!Newtona}% % lang-pl
    Zachodzi % lang-pl
\index{formula!di Newton}% % lang-it
    Vale % lang-it
    \begin{equation}
        \frac{a + b}{c} = \frac{\cos \frac 1 2 (\alpha - \beta)}{\cos \frac 1 2 (\alpha + \beta)}.
    \end{equation}
\end{proposition}

\begin{proposition}
[wzór Mollweidego] % lang-pl
[formula di Mollweide] % lang-it
\index{wzór!Mollweidego}% % lang-pl
\index{formula!di Mollweide}% % lang-it
\label{wzor_mollweidego}%
    Zachodzi % lang-pl
    Vale % lang-it
    \begin{equation}
        \frac{a - b}{c} = \frac{\sin \frac 1 2 (\alpha - \beta)}{\sin \frac 1 2 (\alpha + \beta)}.
    \end{equation}
\end{proposition}

(Nie ma zgodności co do powyższych nazw). % lang-pl
Nieco bardziej geometryczną wersję wzorów poda Izaak Newton w 1707 roku, potem Friedrich von Oppel w 1746. % lang-pl
(Non c'è consenso sui nomi sopra riportati). % lang-it
Una versione leggermente più geometrica delle formule è data da Isaac Newton nel 1707, poi da Friedrich von Oppel nel 1746. % lang-it
\index[persons]{Newton, Izaak}% % lang-pl
\index[persons]{Newton, Isaac}% % lang-it
\index[persons]{Oppel@von Oppel, Friedrich}%
Wyrażenia, których używać będziemy do końca świata i jeden dzień dłużej zawdzięczymy Thomasowi Simpsonowi z 1748 roku, oraz Karlowi Mollweidemu, który opublikuje to samo w 1808 roku bez cytowania poprzedników. % lang-pl
Le espressioni che useremo fino alla fine del mondo e un giorno di più si devono a Thomas Simpson del 1748 e a Karl Mollweide, che pubblicò lo stesso nel 1808 senza citare i predecessori. % lang-it
\index[persons]{Simpson, Thomas}%
\index[persons]{Mollweide, Karl}%

\begin{proposition}
[wzór Regiomontanusa] % lang-pl
[formula di Regiomontano] % lang-it
\index{wzór!Regiomontanusa}% % lang-pl
    Zachodzi % lang-pl
\index{formula!di Regiomontano}% % lang-it
    Vale % lang-it
    \begin{equation}
        \frac{a + b}{a- b} = \frac{\tan \frac 1 2 (\alpha + \beta)}{\tan \frac 1 2 (\alpha - \beta)}.
    \end{equation}
\end{proposition}
% TODO: https://en.wikipedia.org/wiki/Law_of_tangents

Regiomontanus będzie znany gorzej jako Johannes Müller von Königsberg. % lang-pl
Regiomontano è meno noto come Johannes Müller von Königsberg. % lang-it
\index[persons]{Königsberg@von Königsberg, Johannes|see{Regiomontanus}}%
\index[persons]{Regiomontanus}%

\begin{proposition}
    Niech $\alpha, \beta, \gamma$ będą miarami kątów trójkąta o bokach $a, b, c$, polu $S$ i promieniu okręgu opisanego $R$. % lang-pl
    Wtedy % lang-pl
    Siano $\alpha, \beta, \gamma$ le misure degli angoli di un triangolo con lati $a, b, c$, area $S$ e raggio del circolo circoscritto $R$. % lang-it
    Allora % lang-it
    \begin{equation}
        \tan \alpha + \tan \beta + \tan \gamma = \frac{4S}{a^2 + b^2 + c^2 - 8R^2}
    \end{equation}
\end{proposition}

%

% eksperymentalne tłumaczenie na włoski

%

\section{Problem Hansena} % lang-pl
\index{problem!Hansena}% % lang-pl
W trygonometrii problem Hansena to zagadnienie z zakresu pomiarów w terenie, nazwane na cześć astronoma Petra Andreasa Hansena, który będzie pracował przy mapach państwa duńskiego. % lang-pl
\section{Problema di Hansen} % lang-it
\index{problema!di Hansen}% % lang-it
In trigonometria, il problema di Hansen è un problema di topografia, intitolato all'astronomo Peter Andreas Hansen, che lavorò alla realizzazione delle carte del Regno di Danimarca. % lang-it

\begin{problem}
    Dane są dwa znane punkty $A$ i $B$ oraz dwa nieznane punkty $P_1$ i $P_2$. % lang-pl
    Z punktów $P_1$ i $P_2$ obserwator mierzy kąty utworzone przez linie łączące go z każdym z pozostałych trzech punktów skierowane do każdego z pozostałych trzech punktów. % lang-pl
    Wyznaczyć położenie punktów $P_1$ i $P_2$. % lang-pl
    Sono dati due punti noti $A$ e $B$ e due punti ignoti $P_1$ e $P_2$. % lang-it
    Dai punti $P_1$ e $P_2$ un osservatore misura gli angoli formati dalle rette che lo collegano a ciascuno degli altri tre punti. % lang-it
    Determinare la posizione dei punti $P_1$ e $P_2$. % lang-it
\end{problem}

Znamy cztery kąty $\alpha_1 = \angle AP_1B$, $\beta_1 = \angle BP_1 P_2$, $\alpha_2 = \angle  AP_2B$, $\beta_2 = \angle AP_2P_1$. % lang-pl
Wtedy jedno z~rozwiązań zaczyna się od wyznaczenia kątów % lang-pl
Sono noti i quattro angoli $\alpha_1 = \angle AP_1B$, $\beta_1 = \angle BP_1 P_2$, $\alpha_2 = \angle AP_2B$, $\beta_2 = \angle AP_2P_1$. % lang-it
Una delle soluzioni inizia determinando gli angoli % lang-it
\begin{equation}
    \gamma = \pi - \alpha_1 - \beta_1 - \beta_2, \quad
    \delta = \pi - \alpha_2 - \beta_1 - \beta_2
\end{equation}
oraz ilorazu % lang-pl
e il rapporto % lang-it
\begin{equation}
    k = \frac{\sin \gamma \sin \alpha_2 \sin \beta_1}{\sin \delta \sin \alpha_1 \sin \beta_2}.
\end{equation}
Niech % lang-pl
Sia % lang-it
\begin{equation}
    \theta = 2 \arctan \left(\frac{k-1}{k+1} \tan \frac{\beta_1 + \beta_2}{2} \right).
\end{equation}
Wprowadzamy pomocnicze kąty $\varphi = \frac 1 2 (\beta_1 + \beta_2 + \theta)$ oraz $\psi = \frac 1 2 (\beta_1 + \beta_2 - \theta)$. % lang-pl
Szukany odcinek ma długość % lang-pl
Introduciamo gli angoli ausiliari $\varphi = \frac 1 2 (\beta_1 + \beta_2 + \theta)$ e $\psi = \frac 1 2 (\beta_1 + \beta_2 - \theta)$. % lang-it
Il segmento cercato ha lunghezza % lang-it
\begin{equation}
    |P_1 P_2| = |AB| \cdot \frac{\sin \psi \sin \gamma}{\sin \alpha_1 \sin \beta_2} = |AB| \cdot \frac{\sin \varphi \sin \delta}{\sin \alpha_2 \sin \beta_1},
\end{equation}
przy czym by uniknąć błędów numerycznych bierzemy ten ułamek, którego mianownik ma większą wartość bezwzględną. % lang-pl
per evitare errori numerici si utilizza la frazione il cui denominatore ha valore assoluto maggiore. % lang-it

%
% eksperymentalne tłumaczenie na włoski

%

\pol{\section{Problem Snelliusa-Pothenota}}
\ita{\section{Problema di Snellius-Pothenot}}

\pol{Problem Snelliusa-Pothenota rozwiąże po raz pierwszy Willebrord Snellius w 1615 roku.}
\ita{Il problema di Snellius-Pothenot sarà risolto per la prima volta da Willebrord Snellius nel 1615.}
\pol{\index{problem!Snelliusa-Pothenota}}%
\ita{\index{problema!di Snellius-Pothenot}}%

\pol{Wkład Laurenta Pothenota (żyjącego w latach 1650--1732) w problem lub jego rozwiązanie jest niejasny, dlatego George McCaw, słynny brytyjski geodeta stwierdzi, że powinniśmy pomijać jego nazwisko.}
\ita{Il contributo di Laurent Pothenot (1650--1732) al problema o alla sua soluzione non è chiaro, perciò George McCaw, celebre geodeta britannico, sosterrà che dovremmo omettere il suo nome.}

\begin{problem}
    \pol{Dane są trzy znane punkty $A$, $B$ i $C$ oraz nieznany punkt $P$.}
    \ita{Sono dati tre punti noti $A$, $B$ e $C$ e un punto ignoto $P$.}
    \pol{Z punktu $P$ obserwator mierzy kąty utworzone przez linie łączące go z każdym z pozostałych trzech punktów.}
    \ita{Dal punto $P$ un osservatore misura gli angoli formati dalle linee che lo congiungono con ciascuno degli altri tre punti.}
    \pol{Wyznaczyć położenie punktów $P$.}
    \ita{Determinare la posizione del punto $P$.}
\end{problem}

\pol{Zdegenerowany przypadek, kiedy punkt $P$ znajduje się blisko okręgu opisanego na trójkącie $ABC$ powinien być unikany jak ogień, szczególnie jeśli $P$ oznacza położenie rozbitka na morzu, zaś $A$, $B$, $C$ to trzy latarnie na lądzie.}
\ita{Il caso degenere in cui il punto $P$ si trovi vicino al cerchio circoscritto al triangolo $ABC$ dovrebbe essere evitato come la peste, soprattutto se $P$ rappresenta la posizione di un naufrago in mare e $A$, $B$, $C$ sono tre fari sulla costa.}
%


\section{Nierówności trygonometryczne} % lang-pl
\section{Disuguaglianze trigonometriche} % lang-it


\subsection{Nierówność Arystarcha} % lang-pl
\subsection{Disuguaglianza di Aristarco} % lang-it

\begin{theorem}[nierówność Arystarcha] % lang-pl
\begin{theorem}[disuguaglianza di Aristarco] % lang-it
\index{nierówność!Arystarcha}
    Niech $\alpha, \beta$ będą kątami ostrymi takimi, że $\beta < \alpha$. % lang-pl
\index{disuguaglianza!di Aristarco} % lang-it
    
    Wtedy % lang-pl
    Siano $\alpha$ e $\beta$ angoli acuti tali che $\beta < \alpha$. % lang-it
    
    Allora % lang-it
    \begin{equation}
        \frac{\sin \alpha}{\sin \beta} < \frac \alpha \beta < \frac{\tan \alpha}{\tan \beta}.
    \end{equation}
	% https://en.wikipedia.org/wiki/Aristarchus's_inequality
\end{theorem}

Arystarch z Samos będzie greckim astronomem, uczniem Stratona z Lampsaku, który jako pierwszy zaproponuje heliocentryczny model Układu Słonecznego. % lang-pl
Aristarco di Samo fu un astronomo greco, allievo di Stratone di Lampsaco, che per primo propose un modello eliocentrico del Sistema Solare. % lang-it


Stoik Kleantes, następca Zenona z Kition, oskarży go o bezbożność. % lang-pl
Lo stoico Cleante, successore di Zenone di Cizio, lo accusò di empietà. % lang-it

\input{chapters/02-trojkaty/nierownosc-mikolaja_z_kuzy}
\input{chapters/02-trojkaty/nierownosc-snelliusa-huygensa}

%

\section{Zastosowanie: twierdzenie Urquharta} % lang-pl
\section{Applicazione: teorema di Urquhart} % lang-it

\begin{theorem}
[Urquharta] % lang-pl
[di Urquhart] % lang-it
\label{theorem_urquhart}%
\index{twierdzenie!Urquharta}% % lang-pl
\index{teorema!di Urquhart}% % lang-it
    Przy oznaczeniach jak na rysunku: % lang-pl
    Con le notazioni della figura: % lang-it
    \begin{center}
\begin{comment}
    \begin{tikzpicture}[scale=.4]
        \tkzDefPoint(0, 0){E}
        \tkzDefPoint(267:5){E1}
        \tkzDefPoint(250:5){E2}
        \tkzDefPoint(165:5){E3}
        \tkzDefPoint(135:5){E4}

        \tkzDefLine[tangent at=E1](E) \tkzGetPoint{tan1}
        \tkzDefLine[tangent at=E2](E) \tkzGetPoint{tan2}
        \tkzDefLine[tangent at=E3](E) \tkzGetPoint{tan3}
        \tkzDefLine[tangent at=E4](E) \tkzGetPoint{tan4}

        \tkzInterLL(E4,tan4)(E1,tan1) \tkzGetPoint{S1}
        \tkzInterLL(E4,tan4)(E2,tan2) \tkzGetPoint{S2}
        \tkzInterLL(E3,tan3)(E1,tan1) \tkzGetPoint{S3}
        \tkzInterLL(E3,tan3)(E2,tan2) \tkzGetPoint{S4}

        \tkzInterLL(S3,S4)(S1,S2) \tkzGetPoint{Gorny}
        \tkzInterLL(S1,S3)(S2,S4) \tkzGetPoint{Dolny}

        \tkzLabelPoint[below](S1){$A$}
        \tkzLabelPoint[above left](S2){$D$}
        \tkzLabelPoint(S3){$B$}
        \tkzLabelPoint[above right](S4){$C$}
        \tkzLabelPoint(Dolny){$E$}
        \tkzLabelPoint[right](Gorny){$F$}

        \tkzDrawSegments[line width=0.2mm](S1,Dolny S1,Gorny S3,Gorny S2,Dolny)
        \tkzDrawPoints[size=3,color=black,fill=black!50](S1,S2,S3,S4,Dolny,Gorny)
        % konstrukcja z https://en.wikipedia.org/wiki/Ex-tangential_quadrilateral
\end{tikzpicture}
\end{comment}
    \end{center}
    niech $|AB| + |BC| = |AD| + |CD|$. % lang-pl
    Wtedy $|AE| + |CE| = |AF| + |CF|$. % lang-pl
       sia $|AB| + |BC| = |AD| + |CD|$. % lang-it
    Allora $|AE| + |CE| = |AF| + |CF|$. % lang-it
\end{theorem}

Twierdzenie \ref{theorem_urquhart} nie będzie zbyt popularne, napiszą o nim Bogdańska, Neugebauer \cite[s. 97]{neugebauer_2018}, nieznany autor w $\Delta_{84}^{11}$ i pewnie ktoś jeszcze. % lang-pl
Nie do końca wiadomo, kto udowodni je pierwszy: być może będzie to de Morgan w 1841 roku; teza jest też przypadkiem granicznym innego wyniku Chaslesa, znanego około roku 1860. % lang-pl
Wreszcie Dan Pedoe \cite{pedoe_1976} przypisze to twierdzenie Malcolmowi Livingstonowi Urquhartowi\footnote{Matematyk australijski, będzie żyć w latach 1902-1966 i nie opublikuje żadnej pracy.}. % lang-pl

Il teorema \ref{theorem_urquhart} non diventerà particolarmente popolare; ne scriveranno Bogdańska e Neugebauer \cite[s. 97]{neugebauer_2018}, un autore anonimo in $\Delta_{84}^{11}$ e probabilmente anche qualcun altro. % lang-it
Non è del tutto chiaro chi ne fornirà per primo una dimostrazione: forse de Morgan nel 1841; la tesi è anche un caso limite di un altro risultato di Chasles, noto intorno al 1860. % lang-it
Infine Dan Pedoe \cite{pedoe_1976} attribuirà questo teorema a Malcolm Livingston Urquhart\footnote{Matematico australiano, vivrà negli anni 1902--1966 e non pubblicherà alcun lavoro.}. % lang-it

\index[persons]{Pedoe, Daniel}%
\index[persons]{Urquhart, Malcolm Livingston}%

Pompe napisze w $\Delta_{07}^2$ (bez przywołania nazwiska Urquharta) o możliwości dopisania okręgu do czworokąta. % lang-pl
Pompe scriverà in $\Delta_{07}^2$ (senza menzionare il nome di Urquhart) della possibilità di circoscrivere un cerchio a un quadrilatero. % lang-it

%
% eksperymentalne tłumaczenie na włoski

%

\pol{\section{Zastosowanie: punkt i kąt Crelle'a-Brocarda}}
\ita{\section{Applicazione: punto e angolo di Crelle-Brocard}}

\pol{Poznamy teraz dwa wyróżnione punkty trójkąta, opisane w 1875 roku przez oficera francuskiej armii, Henriego Brocarda oraz wiele lat wcześniej przez Augusta Crelle'a \cite{crelle_1816}, założyciela słynnego czasopisma matematycznego.}
\ita{Studieremo ora due punti notevoli del triangolo, descritti nel 1875 dall'ufficiale dell'esercito francese Henri Brocard e molti anni prima da August Crelle \cite{crelle_1816}, fondatore di una celebre rivista matematica.}
\index[persons]{Brocard, Henri}%
\index[persons]{Crelle, August}%

\begin{definition}
\label{punkty_brocarda}%
    \pol{Niech $\triangle ABC$ będzie trójkątem.}
    \ita{Sia $\triangle ABC$ un triangolo.}
    \pol{Punkt $X$ leżący w jego wnętrzu taki, że kąty $\angle XAB$, $\angle XBC$, $\angle XCA$ są równej miary, nazywamy (pierwszym) punktem Crelle'a-Brocarda.}
    \ita{Il punto $X$ interno al triangolo tale che gli angoli $\angle XAB$, $\angle XBC$, $\angle XCA$ abbiano la stessa ampiezza è chiamato (primo) punto di Crelle-Brocard.}
    \pol{\index{punkt!Crelle'a-Brocarda}}%
    \ita{\index{punto!di Crelle-Brocard}}%
    \begin{center}
\begin{comment}
    \begin{tikzpicture}[scale=.75]
        ...
    \end{tikzpicture}
\end{comment}
    \end{center}
    % Construct a circle through vertices A and B, tangent to side BC of the triangle.
    % Symmetrically, construct the other two circles.
    % These three circles intersect at the first Brocard Point of triangle ABC.
    % https://www.geogebra.org/m/MT679Keu
\end{definition}

\pol{Miarę wspomnianych kątów nazywamy kątem Crelle'a-Brocarda.}
\ita{La misura dei suddetti angoli è chiamata angolo di Crelle-Brocard.}
\pol{\index{kąt!Crelle'a-Brocarda}}%
\ita{\index{angolo!di Crelle-Brocard}}%
\pol{Jest jeszcze drugi punkt Crelle'a-Brocarda, gdzie odwracamy kolejność punktów: $\angle XBA = \angle XCB = \angle XAC$.}
\ita{Esiste anche un secondo punto di Crelle-Brocard, ottenuto invertendo l'ordine dei punti: $\angle XBA = \angle XCB = \angle XAC$.}

\begin{proposition}
\label{miara_kata_crellea_brocarda}%
    \pol{W każdym trójkącie istnieje (jedyny) punkt Crelle'a-Brocarda.}
    \ita{In ogni triangolo esiste un unico punto di Crelle-Brocard.}
    \pol{Miara kąta Crelle'a-Brocarda $\omega$ spełnia związek}
    \ita{La misura dell'angolo di Crelle-Brocard $\omega$ soddisfa la relazione}
    \begin{equation}
        \frac 1 {\tan \omega} = \frac 1 {\tan \alpha} + \frac 1 {\tan \beta} + \frac 1 {\tan \gamma},
    \end{equation}
    \pol{gdzie $\alpha, \beta, \gamma$ to miary kątów trójkąta.}
    \ita{dove $\alpha, \beta, \gamma$ sono le ampiezze degli angoli del triangolo.}
    \pol{Ponadto, $0 \le \omega \le \pi/6$.}
    \ita{Inoltre, $0 \le \omega \le \pi/6$.}
\end{proposition}

\pol{Zetel \cite[s. 176]{zetel_2020} pochwali się, że Szebarszin zauważył, że funkcje trygonometryczne kąta Brocarda można obliczyć prosto i że wynikiem będzie}
\ita{Zetel \cite[p. 176]{zetel_2020} osserva che Šebaršin notò come le funzioni trigonometriche dell'angolo di Brocard possano essere calcolate facilmente e che il risultato è}
\index[persons]{Szebarszin, M N}%
\begin{equation}
    \cos = \frac{a^2 + b^2 + c^2}{2\sqrt{a^2b^2 + a^2c^2 + b^2c^2}}.
\end{equation}

\pol{Na 32. Międzynarodowej Olimpiadzie Matematycznej pojawi się następujące zadanie:}
\ita{Alla 32ª Olimpiade Internazionale della Matematica comparve il seguente problema:}

\begin{problem}
    \pol{Niech $P$ będzie dowolnym punktem wewnątrz trójkąta $\triangle ABC$.}
    \ita{Sia $P$ un punto arbitrario all'interno del triangolo $\triangle ABC$.}
    \pol{Wykazać, że co najmniej jeden z kątów $\angle PAB$, $\angle PBC$, $\angle PCA$ jest mniejszy lub równy $\pi /6$.}
    \ita{Dimostrare che almeno uno degli angoli $\angle PAB$, $\angle PBC$, $\angle PCA$ è minore o uguale a $\pi /6$.}
\end{problem}

\pol{Waldemar Pompe (w $\Delta_{92}^{3}$) przedstawi siedem wskazówek, po których zadanie rozwiąże się nieomal samo.}
\ita{Waldemar Pompe (in $\Delta_{92}^{3}$) presenterà sette suggerimenti dopo i quali il problema si risolverà quasi da solo.}
% https://www.deltami.edu.pl/1992/03/krok-po-kroku/

\pol{Bogdańska, Neugebauer \cite[s. 100]{neugebauer_2018} podadzą jako ćwiczenie, (Zetel \cite[s. 148]{zetel_2020} zaś z dowodem):}
\ita{Bogdańska e Neugebauer \cite[p. 100]{neugebauer_2018} lo propongono come esercizio (mentre Zetel \cite[p. 148]{zetel_2020} ne fornisce una dimostrazione):}

\begin{proposition}
    \pol{Niech $X$ będzie punktem Crelle'a-Brocarda trójkąta $\triangle ABC$.}
    \ita{Sia $X$ il punto di Crelle-Brocard del triangolo $\triangle ABC$.}
    \pol{Niech $R_a$, $R_b$ i $R_c$ oznaczają promienie okręgów opisanych na trójkątach $\triangle XBC$, $\triangle AXC$, $\triangle ABX$, zaś $R$ będzie jak zwykle promieniem okręgu opisanego na trójkącie $\triangle ABC$.}
    \ita{Siano $R_a$, $R_b$ e $R_c$ i raggi delle circonferenze circoscritte ai triangoli $\triangle XBC$, $\triangle AXC$, $\triangle ABX$, mentre $R$ sia il raggio della circonferenza circoscritta al triangolo $\triangle ABC$.}
    \pol{Wtedy}
    \ita{Allora}
    \begin{equation}
        R = \sqrt[3]{R_a R_b R_c}.
    \end{equation}
\end{proposition}

\pol{Peter Yiff \cite{yff_1963} postawi w 1963 roku (!) hipotezę, że $8 \omega^3 \le \alpha \beta \gamma$.}
\ita{Peter Yiff \cite{yff_1963} formulò nel 1963 (!) la congettura che $8 \omega^3 \le \alpha \beta \gamma$.}
\index[persons]{Yiff, Peter}%
\pol{Dowód znajdzie w 1974 roku Faruk Abi-Khuzam \cite{abikhuzam_1974}.}
\ita{La dimostrazione sarà trovata nel 1974 da Faruk Abi-Khuzam \cite{abikhuzam_1974}.}
\index[persons]{Abi-Khuzam, Faruk}%

\pol{Mamy jeszcze mało ciekawy dla kogoś o wiedzy tak nikłej jak my okrąg Crelle'a-Brocarda.}
\ita{Abbiamo inoltre la circonferenza di Crelle-Brocard, poco interessante per chi possiede conoscenze limitate come le nostre.}
% PRZECZYTANO: https://mathworld.wolfram.com/BrocardCircle.html
\pol{\index{okrąg!Crelle'a-Brocarda}}%
\ita{\index{circonferenza!di Crelle-Brocard}}%
\pol{Jego średnicą jest odcinek łączący środek okręgu opisanego z punktem Lemoine'a.}
\ita{Il suo diametro è il segmento che unisce il circocentro al punto di Lemoine.}
\pol{\index{okrąg!opisany}}%
\ita{\index{circonferenza!circoscritta}}%
\pol{\index{punkt!Lemoine'a}}%
\ita{\index{punto!di Lemoine}}%
\pol{Przechodzi przez pierwszy i drugi punkt Crelle'a-Brocarda oraz wierzchołki celowo niezdefiniowanego tu trójkąta Brocarda, co uzasadnia stosowaną czasami nazwę ,,okrąg siedmiu punktów''.}
\ita{Passa per il primo e il secondo punto di Crelle-Brocard e per i vertici del triangolo di Brocard, volutamente non definito qui, il che giustifica il nome talvolta usato di \emph{circonferenza dei sette punti}.}
\pol{\index{okrąg!siedmiu punktów}}%
\ita{\index{circonferenza!dei sette punti}}%
\pol{Ma promień równy}
\ita{Ha raggio pari a}
\begin{equation}
    R \cdot \frac{\sqrt{1 - 4 \sin^2 \omega}}{2 \cos \omega}.
\end{equation}

%

% https://en.wikipedia.org/wiki/Rule_of_marteloio - napisać, że tego nie będzie

\section{Bałagan} % lang-pl

\section{Varie} % lang-it

\begin{itemize}
% Twierdzenie Malfattiego. Guzicki-11.
\item Twierdzenie Eulera $1/4R^2$ % lang-pl
\item Twierdzenie Taylora, okrąg, sześciokąt % lang-pl
\item Teorema di Eulero $1/4R^2$ % lang-it
\item Teorema di Taylor, circonferenza, esagono % lang-it
\end{itemize}

%